% =====================================================================
%  lean-stil.tex — Lean 4 kod renklendirme ve kod kutuları
%  Lean 4 ile Modern Mantık
%
%  Renklendirme minted + Pygments ile yapılır (Unicode ve Türkçe
%  karakterler sorunsuz). Kitaba özel sözcük çözümleyicisi:
%  stil/leankitap.py
%
%  Derleme -shell-escape gerektirir:
%      latexmk -xelatex -shell-escape ana
% =====================================================================

\usepackage[newfloat]{minted}
\usemintedstyle{bw}
\tcbuselibrary{minted}

% Sözcük çözümleyicisi: Pygments'e dosyadan özel çözümleyici yükletmek için
% "-x" bayrağı gerekir; dil adı olarak dosya yolu ve sınıf adı verilir.
\renewcommand{\MintedPygmentize}{pygmentize -x}

% ---------- Renk paleti ----------
\definecolor{leanbg}{HTML}{F7F9FC}      % kod arka planı
\definecolor{leanframe}{HTML}{3B5BDB}   % kod çerçevesi (lacivert)
\definecolor{leankw}{HTML}{7048E8}      % anahtar kelime (mor)
\definecolor{leantac}{HTML}{D6336C}     % taktik (pembe-kırmızı)
\definecolor{leantype}{HTML}{0C8599}    % Prop, Type (camgöbeği)
\definecolor{leanhyp}{HTML}{E8590C}     % hipotez adları (turuncu)
\definecolor{leancom}{HTML}{5C7F4A}     % yorum (zeytin yeşili)
\definecolor{leansym}{HTML}{1971C2}     % → ⊢ ∧ ∨ (mavi)
\definecolor{leancmd}{HTML}{862E9C}     % #check, #eval
\definecolor{leankirik}{HTML}{ADB5BD}   % satır numarası

\definecolor{hedefbg}{HTML}{FFF9DB}     % hedef durumu kutusu
\definecolor{hedefframe}{HTML}{F59F00}
\definecolor{bittibg}{HTML}{EBFBEE}
\definecolor{bittiframe}{HTML}{2F9E44}
\definecolor{hatabg}{HTML}{FFF5F5}
\definecolor{hataframe}{HTML}{E03131}

% ---------- Token renkleri ----------
% Pygments'in "bw" stili yüklendikten sonra, her kod bloğunun başında
% token makroları kitabın renk diliyle değiştirilir (formatcom).
\makeatletter
\newcommand{\leanrenk}[3]{% #1: token kısaltması, #2: renk, #3: ek biçim
  \@namedef{PYG@tok@#1}{\def\PYG@tc####1{\textcolor{#2}{####1}}#3}}
\newcommand{\leanrenkleri}{%
  \leanrenk{k}{leankw}{\let\PYG@bf=\textbf}%          anahtar kelime
  \leanrenk{kt}{leantype}{\let\PYG@bf=\textbf}%       Prop, Type
  \leanrenk{nb}{leantac}{\let\PYG@bf=\textbf}%        taktik
  \leanrenk{nv}{leanhyp}{\let\PYG@bf=\textbf}%        hipotez adı
  \leanrenk{nd}{leancmd}{\let\PYG@bf=\textbf}%        #check
  \leanrenk{o}{leansym}{\let\PYG@bf=\relax}%          → ∧ ∨ ⊢
  \leanrenk{nt}{leankw}{\let\PYG@bf=\relax}%          ⟨ ⟩
  \leanrenk{c}{leancom}{\let\PYG@it=\textit}%         yorum
  \leanrenk{c1}{leancom}{\let\PYG@it=\textit}%
  \leanrenk{cm}{leancom}{\let\PYG@it=\textit}%
  \leanrenk{s2}{leantype}{\let\PYG@it=\relax}%        dizgi
  \leanrenk{s}{leantype}{\let\PYG@it=\relax}%
  \leanrenk{nx}{leansym}{\def\PYG@ff####1{{\sembolkodfont ####1}}}% ⋀ ⋁ ⟹ ⊨
}
\makeatother

\setminted{
  fontsize=\small,
  breaklines, breakanywhere,
  tabsize=2,
  numbersep=8pt,
  formatcom=\leanrenkleri,
}
\renewcommand{\theFancyVerbLine}{\textcolor{leankirik}{\tiny\arabic{FancyVerbLine}}}

% Lean dosyalarından çıkarılan kesitlerin yeri (ana.tex'e göre)
\newcommand{\kesitdizini}{kesitler}

% Derleme başında kesitleri yeniden çıkar (shell-escape açıkken)
\AtBeginDocument{%
  \immediate\write18{python3 ../araclar/kesitleri_cikar.py
    ../lean/ModernMantik \kesitdizini}}

% ---------- Kod kutuları ----------
\tcbset{
  leankutu/.style={
    enhanced, breakable,
    colback=leanbg, colframe=leanframe,
    boxrule=0.8pt, arc=3pt,
    left=4pt, right=6pt, top=6pt, bottom=4pt,
    fonttitle=\sffamily\bfseries\small,
    coltitle=white,
    attach boxed title to top left={xshift=8pt,yshift=-2.2mm},
    boxed title style={colback=leanframe, arc=2pt, boxrule=0pt},
    before skip=12pt, after skip=10pt,
    listing engine=minted,
    minted language=stil/leankitap.py:LeanKitapLexer,
    minted options={linenos, xleftmargin=1.6em, fontsize=\small,
                    breaklines, breakanywhere, numbersep=8pt,
                    formatcom=\leanrenkleri},
  }
}

% Kitap içine doğrudan yazılan kod (derlenmeyen örnekler için):
%   \begin{leankod}[Başlık] ... \end{leankod}
\newtcblisting{leankod}[1][Lean 4]{leankutu, title={#1}, listing only}

% Eşlik eden Lean projesinden ÇEKİLEN kod (tek doğruluk kaynağı):
%   \leandosya[Başlık]{Bolum02}{hipotetik_kiyas}
% Bolum02.lean içindeki  -- ANCHOR: hipotetik_kiyas  ...
%                        -- ANCHOR_END: hipotetik_kiyas  arası basılır.
\newcommand{\leandosya}[3][Lean 4]{%
  \tcbinputlisting{leankutu,
    title={#1\quad{\normalfont\scriptsize\ttfamily ModernMantik/#2.lean}},
    listing only,
    listing file={\kesitdizini/#2/#3.lean}}%
}

% Hedef durumu (Lean'in Infoview paneli):
%   \begin{hedefdurumu}[Başlık] ... \end{hedefdurumu}
\newtcblisting{hedefdurumu}[1][Hedef durumu]{
  enhanced, listing only,
  listing engine=minted, minted language=stil/leankitap.py:LeanKitapLexer,
  minted options={fontsize=\small, formatcom=\leanrenkleri},
  colback=hedefbg, colframe=hedefframe,
  boxrule=0.6pt, arc=2pt, leftrule=3pt,
  left=6pt, top=2pt, bottom=2pt,
  title={\sffamily\footnotesize\bfseries #1},
  coltitle=black!75, colbacktitle=hedefframe!30,
  toptitle=1pt, bottomtitle=1pt,
  before skip=6pt, after skip=6pt,
}

% Lean hata mesajı kutusu
\newtcblisting{leanhata}[1][Lean hata mesajı]{
  enhanced, listing only,
  listing engine=minted, minted language=text,
  minted options={fontsize=\small},
  colback=hatabg, colframe=hataframe,
  boxrule=0.6pt, arc=2pt, leftrule=3pt,
  left=6pt, top=2pt, bottom=2pt,
  title={\sffamily\footnotesize\bfseries #1},
  coltitle=hataframe!80!black, colbacktitle=hataframe!15,
  toptitle=1pt, bottomtitle=1pt,
  before skip=6pt, after skip=6pt,
}

% "İspat tamamlandı" çubuğu
\newcommand{\ispatbitti}{%
  \begin{tcolorbox}[enhanced, colback=bittibg, colframe=bittiframe,
      boxrule=0.6pt, arc=2pt, leftrule=3pt,
      left=6pt, top=3pt, bottom=3pt, before skip=6pt, after skip=8pt]
    {\leanfont\small No goals}\hfill
    \textcolor{bittiframe}{\sffamily\bfseries İspat tamamlandı {\sembolfont ✓}}
  \end{tcolorbox}}

% ---------- Metin içi kısayollar ----------
\newcommand{\leantac}[1]{{\leanfont\textcolor{leantac}{\bfseries #1}}}
\newcommand{\leankw}[1]{{\leanfont\textcolor{leankw}{\bfseries #1}}}
\newcommand{\leanhyp}[1]{{\leanfont\textcolor{leanhyp}{\bfseries #1}}}
\newcommand{\leantype}[1]{{\leanfont\textcolor{leantype}{\bfseries #1}}}
\newcommand{\lc}[1]{{\leanfont #1}}

% İki hedef durumunu yan yana gösterme (önce ⇒ sonra):
%   \begin{hedefgecisi}
%     \begin{hedefdurumu}[Önce] ... \end{hedefdurumu}
%   \sonra
%     \begin{hedefdurumu}[Sonra] ... \end{hedefdurumu}
%   \end{hedefgecisi}
\newenvironment{hedefgecisi}
  {\par\smallskip\noindent\begin{minipage}[t]{0.455\linewidth}\vspace{0pt}}
  {\end{minipage}\par\smallskip}
\newcommand{\sonra}{\end{minipage}\hfill
  \begin{minipage}[t]{0.05\linewidth}\vspace{0.9cm}\centering
    \textcolor{leantac}{\Large$\Rightarrow$}\end{minipage}\hfill
  \begin{minipage}[t]{0.455\linewidth}\vspace{0pt}}
